\documentclass[final,pdflatex]{sn-jnl}
%\documentclass[submitted,12pt]{article}
\usepackage[font=footnotesize]{caption}
\usepackage{elsnat}


\newcommand{\TE}{thought experiment\xspace}
\newcommand{\IR}{indeterminacy relations\xspace}
\newcommand{\UR}{uncertainty relations\xspace}
\newcommand{\vW}{von Weizsäcker\xspace}

\title{Measurement Reconstruction vs.\ Trajectory Reconstruction: \VW between Hermann and Popper in the Early Reception of Quantum Mechanics}

\author{Marco Giovanelli}\vspace{0.2cm}\affil{\centering
Università di Torino\\
Department of Philosophy and Educational Sciences\\
Via S. Ottavio, 20 10124 - Torino, Italy\\ \vspace{0.2cm}\texttt{marco.giovanelli@unito.it}\vspace{0.5cm}}
\begin{document}
 
 \abstract{From early on, the \UR were understood to preclude exact predictions of trajectories, while apparently still allowing their retrospective reconstruction with arbitrarily precise position and momentum values. In his Chicago lectures, Heisenberg treated the physical reality of such reconstructed trajectories as a \s{matter of taste}. In 1934, Popper seized on this concession and proposed a \TE that, by using conservation laws, sought to turn a sharp \emph{trajectory reconstruction} into an exact trajectory prediction. Heisenberg and his assistants sharply criticized the proposal. This paper argues that in early 1935 Hermann, drawing on a structurally similar \TE, identified Popper's deeper confusion between context-independent trajectory reconstruction and context-dependent \emph{measurement reconstruction}. It shows that the twenty-two-year-old \vW, then Heisenberg's assistant, played a unique role both in criticizing Popper's arrangement and in providing the thought experiment on which Hermann relied. The paper concludes that, on the eve of the EPR debate, the Hermann--Popper controversy came close to raising the question of whether one may legitimately combine context-dependent attributions of position and momentum into a description of the underlying reality.}

\keywords{Karl Popper, Grete Hermann, history of quantum mechanics, complementarity, thought experiments}


\maketitle

\section*{Acknowledgements}
I am grateful to Andrea Reichenberger for generously sharing unpublished archival material by Grete Hermann.


\section*{Ethical Statement}
\begin{itemize}
\item  Funding: None \item  Conflict of Interest: None declared
\item  Ethical approval: Not required \item  Informed consent: Not required \item  Author contribution: I am the sole author of this paper
\end{itemize}


\end{document}
